\section{Mozgásegyenlet}

\szta
\newpage

\subsection{Dinamika alaptétele}

Mivel kis szögekkel dolgozunk egyszerűsíthetjük a szögfüggvényeket.
\begin{align}
    \sin{\varphi} & = \varphi \\
    \cos{\varphi} & = 1
\end{align}

\begin{equation}
    \sum{M_\text{A}}
    = \theta_\text{A} \cdot \ddot{\varphi}
    = M_\text{c} + M_{{\text{G}}_1} + M_{{\text{G}}_2} + M(t) + M_{\text{k}_1} + M_{\text{k}_2} + M_\text{kt}
\end{equation}

\subsection{Erők}
\begin{align}
     & F_\text{c}
    = c_1(l_3 + l_4) \cdot \dot{\varphi} \cdot \cos{\varphi}
    = c_1(l_3 + l_4) \cdot \dot{\varphi} \\
     & F_{{\text{G}}_1}
    = m_1 \cdot g                        \\
     & F_{{\text{G}}_2}
    = m_2 \cdot g                        \\
     & F_{\text{k}_1}
    = k_1 \cdot l_3 \cdot \sin{\varphi}
    = k_1 \cdot l_3 \cdot \varphi        \\
     & F_{\text{k}_2}
    =k_2 \cdot l_1 \cdot \sin{\varphi}
    =k_2 \cdot l_1 \cdot \varphi
\end{align}

\subsection{Torziós erőrugó előfeszítése}
\staticszta
Ebben a helyzetben egyensúly van, tehát a nyomatékok összege $0$.
\begin{align}
    F_{\text{G}_1} \cdot \frac{l_3 + l_4}{2} - {M_\text{kt}}_0 = 0 \\
    {M_\text{kt}}_0 = \siunit{\Mktzero}{\newton\meter}
\end{align}

\subsection{Nyomatékok}
\begin{align}
     & M_\text{C}
    = -F_\text{c} \cdot (l_3 + l_4)\dot{\varphi} \cdot \cos{\varphi}
    = -F_\text{c} \cdot (l_3 + l_4)\dot{\varphi}              \\
     & M_{{\text{G}}_1}
    = F_{\text{G}_1} \cdot \frac{l_3 + l_4}{2} \cdot \cos{\varphi}
    = F_{\text{G}_1} \cdot \frac{l_3 + l_4}{2}                \\
     & M_{{\text{G}}_2}
    = -F_{\text{G}_2} \cdot \frac{l_1 + l_2}{2} \cdot \sin{\varphi}
    = -F_{\text{G}_2} \cdot \frac{l_1 + l_2}{2} \cdot \varphi \\
     & M(t)
    = M_0 \cdot \sin{\omega t}                                \\
     & M_{\text{k}_1}
    = -F_{\text{k}_1} \cdot l_3 \cdot \cos{\varphi}
    = -F_{\text{k}_1} \cdot l_3                               \\
     & M_{\text{k}_2}
    = -F_{\text{k}_2} \cdot l_1 \cdot \cos{\varphi}
    = -F_{\text{k}_2} \cdot l_1                               \\
     & M_\text{kt}
    = {M_\text{kt}}_0 - k_\text{t} \cdot \varphi
\end{align}

\subsection{Tömegek}
\begin{align}
    m_1 & = \frac{d^2 \cdot \pi}{4} \cdot (l_3 + l_4) \cdot \rho = \siunit{\mone}{\kilogram} \\
    m_2 & = \frac{d^2 \cdot \pi}{4} \cdot (l_1 + l_2) \cdot \rho = \siunit{\mtwo}{\kilogram}
\end{align}

\subsection{Redukált tömeg}
\begin{equation}
    \begin{split}
        \theta_A &= \frac{1}{12} \cdot m_1 \cdot (l_3 + l_4)^2 + m_1 \cdot \left( \frac{l_3 + l_4}{2} \right)^2 \\&+ \frac{1}{12} \cdot m_2 \cdot (l_1 + l_2)^2 + m_2 \cdot \left( \frac{l_1 + l_2}{2} \right)^2 \\[1ex]
        &= \frac{1}{3} \cdot m_1 \cdot (l_3 + l_4)^2 + \frac{1}{3} \cdot m_2 \cdot (l_1 + l_2)^2 = \siunit{\thetaA}{\kilogram\meter^2}
    \end{split}
\end{equation}

\newpage
\subsection{Összegzés}
\begin{align}
    \begin{split}
        \theta_{\text{A}} \ddot{\varphi} ={} & -c_1(l_3+l_4)^2 \dot{\varphi} - m_2 g \frac{l_1+l_2}{2} \varphi \\
        & + M_0 \sin(\omega t) - k_1 l_3^2 \varphi - k_2 l_1^2 \varphi - k_t \varphi
    \end{split}                                                                      \\[2ex]
    \ddot{\varphi} + \frac{c_1(l_3+l_4)^2}{\theta_{\text{A}}} \dot{\varphi} & + \frac{\left(k_1 l_3^2 + k_2 l_1^2 + k_{\text{t}} + m_2 g \frac{l_1+l_2}{2}\right)}{\theta_{\text{A}}} \varphi
    = \frac{M_0 \sin(\omega t)}{\theta_{\text{A}}}
\end{align}

A referencia alakból ezek a kifejezések adódnak.
\begin{align}
    \omega_n^2           & = \frac{k_1 \cdot l_3^2 + k_2 \cdot l_1^2 + k_t + m_2 \cdot g \cdot \frac{l_1 + l_2}{2}}{\theta_A} \\[1ex]
    2\zeta\omega_n       & = \frac{c_1(l_3 + l_4)^2}{\theta_A}                                                                \\[1ex]
    f_0 \cdot \omega_n^2 & = \frac{M_0}{\theta_A}
\end{align}